\documentclass[11pt]{article}
\usepackage{amsmath, amssymb}
\usepackage[margin=1in]{geometry}

\title{Algebraic Ideal Vapor-Compression Cycle for R32:\\
Governing Equations (No Simultaneous Solver)}
\author{Shilpa Narasimhan}
\date{\today}

\begin{document}
\maketitle

\section{Why an Equation of State Is Still Required}

Removing the simultaneous nonlinear solver (IPOPT, as used by IDAES) does not
remove the need for an equation of state (EoS). The solver and the EoS serve
different roles. IDAES uses IPOPT to solve the full flowsheet --- every unit
model's mass and energy balances, phase-equilibrium constraints, and bounds
--- as one large coupled nonlinear system. That simultaneous solve is what
introduces path-dependent convergence behavior, a tolerance floor
($\text{tol}=10^{-4}$), and warmstart-sensitive failures unrelated to the
underlying physics.

The equation of state, by contrast, is the physical model itself: it is what
allows real (non-ideal) enthalpy and entropy to be computed near the
saturation dome, where the refrigerant spends the entire cycle. An
ideal-gas heat capacity alone cannot describe this region. Evaluating the
Peng--Robinson EoS at a known temperature and pressure requires no iterative
solver: the cubic in compressibility factor $Z$ is solved directly via its
closed-form/eigenvalue root (deterministic, no initial guess, no failure
mode), and the departure functions built on top of it are closed-form
integrals. The only root-finding anywhere in this formulation is a single
well-posed one-dimensional bisection (Sec.~\ref{sec:state2}), fundamentally
different in kind from a large coupled simultaneous solve.

\section{Peng--Robinson Equation of State}

\begin{equation}
P = \frac{RT}{V-b} - \frac{a(T)}{V(V+b) + b(V-b)}
\end{equation}

with temperature-dependent attraction parameter

\begin{equation}
a(T) = 0.45724\,\frac{R^2 T_c^2}{P_c}\,\alpha(T), \qquad
\alpha(T) = \left[1 + \kappa\left(1 - \sqrt{T/T_c}\right)\right]^2
\end{equation}

\begin{equation}
\kappa = 0.37464 + 1.54226\,\omega - 0.26992\,\omega^2
\end{equation}

and size parameter

\begin{equation}
b = 0.07780\,\frac{R\,T_c}{P_c}
\end{equation}

where $T_c$, $P_c$ are the critical temperature and pressure and $\omega$ is
the acentric factor. In compressibility-factor form, with
$A = a(T)P/(RT)^2$ and $B = bP/(RT)$:

\begin{equation}
Z^3 - (1-B)Z^2 + \left(A - 2B - 3B^2\right)Z - \left(AB - B^2 - B^3\right) = 0
\end{equation}

This cubic is solved directly for its real roots; the largest admissible
root ($Z > B$) is the vapor phase, the smallest is the liquid phase.

\section{Departure Functions (Poling, Table 6-3; PR: $u=2$, $w=-1$)}

Real-fluid enthalpy and entropy are expressed as an ideal-gas baseline plus
a departure term computed from the EoS:

\begin{equation}
h(T,P) = h^{\text{ig}}(T) + h^{\text{dep}}(T,P,Z), \qquad
s(T,P) = s^{\text{ig}}(T,P) + s^{\text{dep}}(T,P,Z)
\end{equation}

\begin{equation}
h^{\text{dep}} = RT(Z-1) + \frac{T\dfrac{da}{dT} - a}{2\sqrt{2}\,b}
\ln\!\left(\frac{Z+(1+\sqrt{2})B}{Z+(1-\sqrt{2})B}\right)
\end{equation}

\begin{equation}
s^{\text{dep}} = R\ln(Z-B) + \frac{\dfrac{da}{dT}}{2\sqrt{2}\,b}
\ln\!\left(\frac{Z+(1+\sqrt{2})B}{Z+(1-\sqrt{2})B}\right)
\end{equation}

\section{Ideal-Gas Enthalpy and Entropy (Shomate)}

With $t = T/1000$ and Shomate coefficients $A$--$E$:

\begin{equation}
h^{\text{ig}}(T) = 1000\left(At + \frac{Bt^2}{2} + \frac{Ct^3}{3}
+ \frac{Dt^4}{4} - \frac{E}{t}\right)
\end{equation}

\begin{equation}
s^{\text{ig}}(T,P) = A\ln t + Bt + \frac{Ct^2}{2} + \frac{Dt^3}{3}
- \frac{E}{2t^2} - R\ln\!\left(\frac{P}{P_{\text{ref}}}\right)
\end{equation}

\section{Saturation Pressure (Ambrose--Walton, Poling \S7-4)}

With reduced temperature $T_r = T/T_c$ and $\tau = 1 - T_r$:

\begin{align}
f^{(0)} &= \frac{-5.97616\,\tau + 1.29874\,\tau^{1.5} - 0.60394\,\tau^{2.5}
- 1.06841\,\tau^5}{T_r} \\
f^{(1)} &= \frac{-5.03365\,\tau + 1.11505\,\tau^{1.5} - 5.41217\,\tau^{2.5}
- 7.46628\,\tau^5}{T_r} \\
f^{(2)} &= \frac{-0.64771\,\tau + 2.41539\,\tau^{1.5} - 4.26979\,\tau^{2.5}
+ 3.25259\,\tau^5}{T_r}
\end{align}

\begin{equation}
P_{\text{sat}}(T) = P_c \exp\!\left(f^{(0)} + \omega f^{(1)} + \omega^2 f^{(2)}\right)
\end{equation}

This is evaluated directly and requires no iteration; it is what fixes
$P_{\text{evap}}$ and $P_{\text{cond}}$ once the saturation temperatures are
specified.

\section{The Cycle: Sequential, Not Simultaneous}

Operating specification: evaporator saturation temperature $T_{\text{evap}}$
fixed; condenser saturation temperature $T_{\text{cond}} = T_{\text{amb}} +
\Delta T_{\text{approach}}$; superheat $= $ subcool $= 0$; compressor
isentropic efficiency $\eta$. Because the saturation temperatures are fixed
and there is no superheat or subcooling, each state can be evaluated from
information already known before it --- the system is sequential rather
than a coupled simultaneous system.

\subsection*{State 1 --- Evaporator Outlet (Saturated Vapor)}

\begin{equation}
P_{\text{evap}} = P_{\text{sat}}(T_{\text{evap}}), \qquad
h_1 = h_{\text{vap}}(T_{\text{evap}}, P_{\text{evap}}), \qquad
s_1 = s_{\text{vap}}(T_{\text{evap}}, P_{\text{evap}})
\end{equation}

Fully determined; no iteration.

\subsection*{State 2 --- Compressor Outlet}
\label{sec:state2}

The isentropic discharge temperature $T_{2s}$ is the unique root of

\begin{equation}
s_{\text{vap}}(T_{2s}, P_{\text{cond}}) - s_1 = 0
\end{equation}

Since $s_{\text{vap}}(T, P_{\text{cond}})$ is monotonically increasing in
$T$, this is a well-posed one-dimensional bisection (bracketed root find,
e.g.\ Brent's method) --- the only implicit step in the entire cycle, and
categorically different from a coupled multi-equation solve: one variable,
one guaranteed-monotonic function, no path dependence.

\begin{equation}
h_{2s} = h_{\text{vap}}(T_{2s}, P_{\text{cond}}), \qquad
h_2 = h_1 + \frac{h_{2s} - h_1}{\eta}
\end{equation}

\subsection*{State 3 --- Condenser Outlet (Saturated Liquid)}

\begin{equation}
P_{\text{cond}} = P_{\text{sat}}(T_{\text{cond}}), \qquad
h_3 = h_{\text{liq}}(T_{\text{cond}}, P_{\text{cond}}), \qquad
s_3 = s_{\text{liq}}(T_{\text{cond}}, P_{\text{cond}})
\end{equation}

Fully determined; no iteration.

\subsection*{State 4 --- Valve Outlet (Isenthalpic Expansion)}

\begin{equation}
h_4 = h_3
\end{equation}

a direct assignment, not a calculation. The resulting state is two-phase at
$P_{\text{evap}}$; its vapor quality follows by linear interpolation between
the saturated-liquid and saturated-vapor enthalpies at $T_{\text{evap}}$
(already available from State 1 and the corresponding liquid evaluation),
again with no root-finding required:

\begin{equation}
x_4 = \frac{h_4 - h_{f,4}}{h_{g,4} - h_{f,4}}, \qquad
s_4 = s_{f,4} + x_4\left(s_{g,4} - s_{f,4}\right)
\end{equation}

\section{Coefficient of Performance}

\begin{equation}
q_{\text{evap}} = h_1 - h_4, \qquad w_{\text{comp}} = h_2 - h_1
\end{equation}

\begin{equation}
\boxed{\text{COP} = \dfrac{q_{\text{evap}}}{w_{\text{comp}}}
= \dfrac{h_1 - h_4}{h_2 - h_1}}
\end{equation}

Because COP is a ratio of enthalpy \emph{differences}, any constant
reference-state offset in $h^{\text{ig}}$ or $s^{\text{ig}}$ cancels
exactly; no reference-state anchoring is required.

\section{Summary of the Solve Structure}

\begin{center}
\begin{tabular}{lll}
\textbf{Step} & \textbf{Type} & \textbf{Method} \\
\hline
$P_{\text{sat}}$ (states 1, 3) & closed-form & Ambrose--Walton \\
$Z$ (each state) & closed-form & cubic root (direct) \\
$h,\,s$ (each state) & closed-form & Shomate $+$ PR departure \\
State 2 discharge $T$ & 1-D root find & bracketed bisection (monotone) \\
State 4 quality & closed-form & linear interpolation \\
\hline
Full flowsheet (IDAES/IPOPT) & simultaneous solve & \emph{not used here} \\
\end{tabular}
\end{center}

Every deviation case (Tc, Pc perturbed) is therefore evaluated
deterministically: one answer per case, no convergence flag, no warmstart
dependence, and no tolerance floor obscuring small perturbations.

\end{document}
